% Chapter 1: Generalities. Background needed to read Chapters 2 and 3, and the place of
% BLOC-BOT among existing robots. Sources: references.bib (each claim cited).

\chapter{Generalities}
\label{ch:generalities}

\section{Introduction}
This chapter gives the background on which the rest of the thesis builds. It first
describes mobile robots and the way they move, then looks at robots designed for teaching
and at the idea of a modular robot, and places BLOC-BOT among them. The last two sections
introduce the electronic building blocks of the project: microcontrollers, and the wireless
links that let two of them communicate.

\section{Mobile robots}
\label{sec:mobile-robots}

\subsection{Definition}
A mobile robot is a machine that can move through its environment under its own control,
as opposed to an industrial arm fixed to the floor of a factory. It carries its own
actuators to move, sensors to perceive its surroundings and its own motion, a processing
unit that decides what to do, and usually its own energy source. It may be controlled by a
human, as BLOC-BOT is when it is driven from its controller, or it may act on its own from
what its sensors tell it, as a line-following robot does.

\subsection{Locomotion}
A mobile robot can move on legs, on wheels, on tracks, or in more unusual ways. Legged
locomotion copes well with soft and uneven ground but needs many degrees of freedom and is
mechanically complex. Wheels, by contrast, are simple and very well suited to flat ground,
where wheeled locomotion is one to two orders of magnitude more efficient than legged
locomotion \cite{siegwart2004}. For these reasons, Siegwart and Nourbakhsh describe the
wheel as ``by far the most popular locomotion mechanism in mobile robotics''
\cite{siegwart2004}. The student competitions the authors took part in are run on flat indoor
tracks, where wheels are the natural choice.

\subsection{Differential drive}
\label{sec:diff-drive}
Among the many possible wheel arrangements \cite{siegwart2004}, the simplest steerable
one is the \emph{differential drive}: two wheels on a common axis, each driven by its own
motor, and one or two passive points of contact, such as a caster wheel, that keep the
robot level. There is no steering mechanism. The robot goes straight when both wheels turn
at the same speed, turns when one wheel turns faster than the other, and spins on the spot
when the wheels turn at the same speed in opposite directions. This ability to turn
without moving its centre makes it very manoeuvrable, and Siegwart and Nourbakhsh name the
two-wheel differential drive the most popular of these highly manoeuvrable designs
\cite{siegwart2004}.

The same property has a cost. Since nothing mechanical keeps the robot on a straight line,
both motors ``must be driven along exactly the same velocity profile'' for it to go
straight, which is difficult because of the small differences between wheels and motors
\cite{siegwart2004}. This is why a differential-drive robot
benefits from measuring the speed of each wheel, as BLOC-BOT does with its encoders
(Section~\ref{sec:encoders}). BLOC-BOT uses this configuration: two driven wheels and a
caster at the front.

\section{Educational and competition robots}
\label{sec:edu-robots}
Robots are widely used to teach electronics, programming and control, because they turn
abstract notions into a machine whose behaviour can be seen. Several families of robots
serve this purpose.

\paragraph{Robots designed for teaching.} Some robots are designed from the start as
teaching tools. The e-puck, developed at the EPFL in Switzerland, is a small
differential-drive robot made for engineering education at university level, inexpensive
enough to be used widely in class, released under an open-hardware licence, and extensible
with additional boards stacked on top of it, for example for radio communication
\cite{mondada2009epuck}. Thymio, also from the EPFL, targets a wider public, from primary
school to university: it is open-source in both hardware and software, is driven by two
independent motors, and can be used through pre-programmed behaviours, then graphical and
text programming \cite{riedo2013thymio}. Both are complete, finished robots: they are
excellent to learn with, but their mechanics are fixed, and adding a sensor that was not
planned by their designers is not their purpose.

\paragraph{Hobby kits.} At the other end are the robot-car kits sold in electronics shops.
The most common one in Algerian shops, a two-wheel chassis kit, costs 2950~DA and contains
a chassis, two 6~V gearmotors with a 1:48 reduction, two wheels, an omnidirectional caster,
two encoder discs, a holder for four AA batteries and screws; the microcontroller, the
motor driver and every sensor have to be bought separately \cite{mhtronic_kit}. Such kits
are cheap and are the usual starting point for student teams, but they provide little more
than a frame and motors. Every new sensor needs its own way of being fixed and wired, and
the result is a different robot every time.

\paragraph{Competition robots.} In the authors' experience of competitions such as the Algerian
Robot Cup (Section~\ref{sec:spec}), teams build their robots from such kits or from
scratch, for one event and one task. As explained in Chapter~\ref{ch:design}, this is precisely the
situation BLOC-BOT was created to improve.

\section{Modular robots}
\label{sec:modular}
The word \emph{modular} is used in robotics with two rather different meanings, and it is
worth separating them to place BLOC-BOT correctly.

In research, \emph{modular self-reconfigurable robots} are machines made of many identical
or similar modules that can rearrange how they are connected, and so change their own
shape, to adapt to a new situation, perform a new task or recover from damage
\cite{yim2007modular}. The robot changes its shape by itself. This is an active research field
with its own open problems \cite{yim2007modular}.

In a simpler sense, a robot is modular when it is built from parts that follow a common
interface, so that they can be added, removed or exchanged by hand without redesigning the
rest. The extension boards of the e-puck are an example of this on the electronic side
\cite{mondada2009epuck}.

BLOC-BOT belongs to this second kind. It does not reconfigure itself: its modularity is
mechanical and manual. A plate with a fixed pattern of holes, the waffle plate, provides
a common interface, and printed holders, the blocs, clamp a sensor or a module to it at
any position of the pattern. Changing the robot for a new competition then means moving or
adding blocs and their wiring, rather than designing and printing a new chassis. In this
sense, BLOC-BOT sits between the hobby kit, which it resembles in cost and in parts, and the
educational robot, from which it takes the idea of a documented, open platform that others
can build on.

\section{Embedded systems and microcontrollers}
\label{sec:mcu}
An embedded system is a computer built into a device to perform a specific function, often
with tight limits on size, cost and power, and directly connected to the physical world
through sensors and actuators. A robot's brain is usually such a system, built around a
\emph{microcontroller}: a single chip that contains a processor, program and data memory,
and the peripherals needed to talk to the outside world, such as general-purpose
input/output (GPIO) pins, timers that generate PWM signals, analog-to-digital converters
(ADCs) and serial interfaces.

\paragraph{Arduino.} Much of student robotics starts with the Arduino UNO, a board built
around the ATmega328P microcontroller \cite{arduino_uno}. Its 8-bit processor runs at
16~MHz, with 32~KB of flash memory for the program, 2~KB of SRAM and a 10-bit ADC
\cite{atmega328p_datasheet}. The board has no radio \cite{arduino_uno}: a wireless link
needs an extra module.

\paragraph{The ESP32 family.} The ESP32 chips from Espressif are microcontrollers with an
integrated 2.4~GHz radio. The original ESP32 combines a dual-core 32-bit processor at up to
240~MHz, 520~KB of SRAM, Wi-Fi and Bluetooth~4.2 \cite{esp32_datasheet}; the newer ESP32-S3
offers a similar processor with 512~KB of SRAM, Wi-Fi and Bluetooth~5 Low Energy, and more
GPIOs \cite{esp32s3_datasheet}. Table~\ref{tab:mcu-compare} compares them with the Arduino
UNO's microcontroller. Their main interest for a robot is that the processor and the radio
are on the same chip, so the robot and its controller can communicate without any added
module, and that they are programmed with the same Arduino tools as an Arduino board.

\begin{table}[htbp]
  \centering
  \caption[Microcontrollers compared]{Microcontrollers compared
           \cite{atmega328p_datasheet,arduino_uno,esp32_datasheet,esp32s3_datasheet}.}
  \label{tab:mcu-compare}
  \begin{tabularx}{\textwidth}{@{}l >{\raggedright\arraybackslash}X >{\raggedright\arraybackslash}X >{\raggedright\arraybackslash}X@{}}
    \toprule
    & ATmega328P (Arduino UNO) & ESP32 & ESP32-S3 \\
    \midrule
    Processor & 8-bit AVR, 16~MHz & dual-core 32-bit Xtensa LX6, up to 240~MHz
              & dual-core 32-bit Xtensa LX7, up to 240~MHz \\
    SRAM & 2~KB & 520~KB & 512~KB \\
    Radio & none & Wi-Fi 802.11b/g/n, Bluetooth~4.2 (Classic and LE)
          & Wi-Fi 802.11b/g/n, Bluetooth~5 (LE only) \\
    ADC & 10-bit & two 12-bit & two 12-bit \\
    Used in BLOC-BOT & no & controller & robot \\
    \bottomrule
  \end{tabularx}
\end{table}

\section{Wireless links for robots}
\label{sec:wireless}
A robot driven by hand needs a link between its controller and its own microcontroller.
Most of the options available to a student project work in the 2.4~GHz band.

\begin{itemize}
  \item \textbf{Wi-Fi} (IEEE 802.11b/g/n) is built into ESP32 chips \cite{esp32_datasheet}.
        It is designed to connect devices to a network: one device acts as an access
        point and the others join it before exchanging data, usually over IP.
  \item \textbf{Bluetooth} exists in two forms. Bluetooth Classic is available on the
        original ESP32 only; Bluetooth Low Energy,
        available on both chips \cite{esp32_datasheet,esp32s3_datasheet}, organises data as
        services and characteristics exchanged over a connection between a central and a
        peripheral device.
  \item \textbf{Separate radio modules}, such as the nRF24L01+, are transceivers that a
        microcontroller without a radio drives over SPI. The nRF24L01+ works in the 2.4~GHz
        band at 250~kbit/s, 1 or 2~Mbit/s \cite{nrf24l01p_datasheet}; it is cheap and simple
        but adds a module and its wiring at each end.
  \item \textbf{ESP-NOW} is Espressif's connectionless protocol, which uses the Wi-Fi radio
        of ESP32 chips to send short frames directly from one device to another, without an
        access point or a connection set-up, with up to 250~bytes per frame
        \cite{espnow_idf}.
\end{itemize}

The choice between them for BLOC-BOT is made in Section~\ref{sec:espnow}: since both ends
are ESP32 chips, ESP-NOW gives a direct link with no extra hardware and no connection to
establish.

\section{Conclusion}
A small differential-drive robot is one of the simplest and most widely used forms of
wheeled mobile robot \cite{siegwart2004}; its weakness, keeping a straight line, is what wheel
encoders help with. Between finished educational robots, whose mechanics are fixed, and
hobby kits, which offer little beyond a frame and motors, there is room for a platform that
is as cheap as a kit but provides a standard way of mounting and wiring modules. The ESP32
family, with its integrated radio and ESP-NOW, provides the electronic base for such a
platform. Chapter~\ref{ch:design} specifies and designs it.
